% !PRINTED 38
\weberpage{38}
%<<CONT
wie man, wenn $S_{\nu-1}$ noch nicht die ganze Gruppe $S$ erschöpft, daraus eine eben solche umfassendere Gruppe $S_\nu$ ableiten kann.

Für jedes Element $\Theta$ von $S$ wird es einen gewissen niedrigsten positiven Exponenten $a_\nu$ geben, für den $\Theta^{a_\nu}$ in $S_{\nu-1}$ enthalten ist, und dies $a_\nu$ ist wegen (16) gleich oder kleiner als $a_{\nu-1}$, und ist ausserdem eine Potenz von $p$, da es ein Theiler des Grades von $\Theta$ sein muss. Wir wählen $\Theta$ so, dass der Exponent $a_\nu$ so gross als möglich wird. Ist $\Theta_1$ ein beliebiges anderes Element in $S$, und $\Theta_1^h$ die niedrigste Potenz von $\Theta_1$, die in $S_{\nu-1}$ enthalten ist, so ist auch $h$ eine Potenz von $p$ und gleich oder kleiner als $a_\nu$. Es ist daher $\Theta_1^{a_\nu}$ in $S_{\nu-1}$ enthalten.

Setzen wir aber
\weq{17}{\Theta_1^{a_\nu} = A_1^{\lambda_1} A_2^{\lambda_2} \ldots A_{\nu-1}^{\lambda_{\nu-1}},}
so sind die sämmtlichen Exponenten $\lambda$ durch $a_\nu$ theilbar. Denn es ist
\[
\Theta_1^{a_{\nu-1}} = A_1^{\frac{\lambda_1 a_{\nu-1}}{a_\nu}} A_2^{\frac{\lambda_2 a_{\nu-1}}{a_\nu}} \ldots A_{\nu-1}^{\frac{\lambda_{\nu-1} a_{\nu-1}}{a_\nu}},
\]
und nach der Voraussetzung 2) müssen die Exponenten von $A_1, A_2 \ldots A_{\nu-1}$ auf der rechten Seite dieser Formel durch $a_{\nu-1}$ theilbare ganze Zahlen sein. Folglich sind $\lambda_1, \lambda_2 \ldots \lambda_{\nu-1}$ durch $a_\nu$ theilbar. Wir können also, indem wir zu dem oben betrachteten speciellen $\Theta$ zurückkehren und unter $h_1, h_2 \ldots h_{\nu-1}$ ganze Zahlen verstehen,
\[
\Theta^{a_\nu} = A_1^{h_1 a_\nu} A_2^{h_2 a_\nu} \ldots A_{\nu-1}^{h_{\nu-1} a_\nu}
\]
setzen, und wenn wir dann
\[
A_\nu = \Theta A_1^{-h_1} A_2^{-h_2} \ldots A_{\nu-1}^{-h_{\nu-1}}
\]
annehmen, so ist
\weq{18}{A_\nu^{a_\nu} = 1}
die niedrigste Potenz von $A_\nu$, die in $S_{\nu-1}$ enthalten ist (weil sonst auch noch eine niedrigere Potenz von $\Theta$ in $S_{\nu-1}$ enthalten wäre).

Dann ist
\weq{19}{A_1^{\alpha_1} A_2^{\alpha_2} \ldots A_\nu^{\alpha_\nu}, \quad \alpha_i = 0, 1, 2 \ldots a_i - 1}
nur $= 1$, wenn $\alpha_1, \alpha_2 \ldots \alpha_\nu$ durch $a_1, a_2 \ldots a_\nu$ theilbar und folglich $= 0$ sind, und demnach sind die in (19) dargestellten Elemente alle von einander verschieden. Diese Elemente bilden aber eine Gruppe $S_\nu$ mit der Basis $A_1, A_2 \ldots A_\nu$, die der Forderung 1) genügt.
